\section*{Hoofdstuk \ref{ch:b2:dipole-radiation} --- Dipoolstraling en verstrooiing}

\begin{solution}{exo:b2:dipole-radiation:1}
$p_0 = I_0\ell/\omega = \qty{3.2e-9}{C.m}$;
$E = \mu_0\omega^2p_0\sin\theta/4\pi r$: \qty{0.13}{V/m} bij
$\theta = 90^\circ$, \qty{0.063}{V/m} bij \qty{30}{\degree};
$\mathcal P = \mu_0\omega^4p_0^2/12\pi c =
\qty{175}{W}$;
$R_{\text{str}} = 2\mathcal P/I_0^2 = \qty{88}{\ohm}$.
\end{solution}

\begin{solution}{exo:b2:dipole-radiation:2}
$790(\ell/\lambda)^2$: \qty{8.8e-5}{\ohm}, \qty{8.8e-3}{\ohm},
\qty{0.88}{\ohm}, \qty{88}{\ohm}. Koper: $0.026$, $0.084$, $0.27$,
\qty{0.84}{\ohm} (\omterm{thm:b2:waves-in-media:skin}{huideffect}): rendementen $0.3\%$, $10\%$, $77\%$, $99\%$.
Bij lange golflengten heeft alleen een mast die een behoorlijke fractie van
$\lambda$ is, een \omterm{ex:b2:dipole-radiation:antenna}{stralingsweerstand} boven zijn verliezen.
\end{solution}

\begin{solution}{exo:b2:dipole-radiation:3}
$400 : 550 : 700 = 9.4 : 2.6 : 1$. $(10/3)^4 = 120$: de radar van
\qty{3}{cm} ziet kleine druppels honderd keer beter; die van \qty{10}{cm}
kijkt door de regen heen naar wat erachter zit.
\end{solution}

\begin{solution}{exo:b2:dipole-radiation:4}
$\sin^2\chi/(1 + \cos^2\chi)$: $0.14$, $0.60$, $1$, $0.60$. Richt op
\qty{90}{\degree} van de zon (een grote cirkel over de hemel); staat de zon
in het zenit, dan ligt de hele horizon op \qty{90}{\degree} en is de
polarisatie erlangs gericht. Meervoudige verstrooiing, de weerkaatsing van
de grond en de anisotropie van de moleculen voegen ongepolariseerd licht
toe: hoogstens ongeveer $75\%$.
\end{solution}

\begin{solution}{exo:b2:dipole-radiation:5}
(a) $\int\sin^2\theta\,\dd\Omega = 8\pi/3$. (b)
$\int_{60^\circ}^{120^\circ}\sin^3\theta\,\dd\theta/
(4/3) = 0.917/1.333 = 69\%$.
(c) $\sin^2\theta = \tfrac12$: $\theta = \qty{45}{\degree}$, een bundel van
\qty{90}{\degree} breed. (d) Maximum $1$ tegen het gemiddelde $2/3$: $1.5$.
\end{solution}

\begin{solution}{exo:b2:dipole-radiation:6}
(a) $p/4\pi\varepsilon_0r^3 = \omega^2p/4\pi\varepsilon_0c^2r$ bij
$r_0 = c/\omega = \lambda/2\pi$. (b) \qty{48}{m}: op \qty{10}{m} het
quasistatische veld, op \qty{10}{km} het uitgestraalde. (c)
$r_0 = \qty{4.8}{cm}$: het hoofd zit in de nabijheidszone. (d) In de
nabijheidszone zijn $\vect E$ en $\vect B$ in kwadratuur: de \omterm{thm:b2:poynting-vector:poynting}{poyntingvector}
is gemiddeld nul, energie wordt opgeslagen en teruggegeven, zoals in een
condensator.
\end{solution}

\begin{solution}{exo:b2:dipole-radiation:7}
(a) $8\pi r_e^2/3 = \qty{6.65e-29}{m^2}$. (b)
$n_e\sigma_TL = 10^{14} \times 6.65 \times 10^{-29}
\times 10^9 = 7 \times 10^{-6}$:
een miljoenste van de fotosfeer, kleurloos omdat \omterm{prop:b2:dipole-radiation:rayleigh}{thomsonverstrooiing} dat is,
en verdronken in het verstrooide zonlicht van de hemel behalve tijdens een
eclips. (c) \qty{10}{keV} $\gg$ \qty{300}{eV}: de elektronen reageren alsof
zij vrij zijn, $6\sigma_T = \qty{4e-28}{m^2}$. (d) Foto-elektrische
absorptie groeit als $Z^4$ en licht het calcium eruit; verstrooiing is per
elektron in bot en vlees hetzelfde.
\end{solution}

\begin{solution}{exo:b2:dipole-radiation:8}
(a) $\tau = 0.36$, $0.10$, $0.038$: $T = 0.70$, $0.90$, $0.96$. (b)
$\tau \times 38 = 13.6$, $3.8$, $1.45$: $T = 1 \times 10^{-6}$, $0.02$,
$0.23$; rood gedeeld door blauw $\sim 2 \times 10^5$. (c) De dampkring van
de aarde buigt het licht van de zonsondergang de schaduw in, verrood door
dezelfde verstrooiing. (d) Bij de horizon is de weg lang: het blauw wordt
uit het verstrooide licht zelf weggestrooid, en meervoudige verstrooiing
voegt wit toe.
\end{solution}

\begin{solution}{exo:b2:dipole-radiation:9}
(a) Langs de lijn van het paar geeft het wegverschil $\lambda/2$ uitdoving;
dwars erop optelling. (b) In tegenfase: het omgekeerde --- straling langs de
lijn (langsstraler), niets dwars. (c) Naar $+x$: een looptijdfase $-\pi/2$
plus een voedingsfase $+\pi/2$, dus in fase; naar $-x$: $-\pi$, dus
uitdoving: een cardioïde, eenzijdig. (d) Masten vormen het dekkingsgebied en
beschermen buurzenders; een fasegestuurde rij richt zich door de
voedingsfasen te veranderen, in microseconden.
\end{solution}

\begin{solution}{exo:b2:dipole-radiation:10}
(a) $\langle\mathcal P\rangle = e^2\omega_0^4x_0^2/12\pi\varepsilon_0c^3$,
$W = \tfrac12m\omega_0^2x_0^2$:
$\Gamma = \mathcal P/W = e^2\omega_0^2/6\pi\varepsilon_0mc^3$. (b)
$\omega_0 = \qty{3.2e15}{rad/s}$: $\Gamma =
\qty{6.4e7}{s^{-1}}$ --- de
gebruikte waarde. (c) \qty{16}{ns};
$\Delta\lambda = \lambda^2\Gamma/2\pi c =
\qty{1.2e-14}{m}$, een honderdste
picometer (\qty{10}{MHz}). (d) $Q = 5 \times 10^7$.
\end{solution}

\begin{solution}{exo:b2:dipole-radiation:11}
(a) \qty{75}{m}; $I_0 = \sqrt{2\mathcal P/R} = \qty{53}{A}$. (b)
$\langle\Pi\rangle = 1.5 \times
5 \times 10^4/2\pi \times 10^8 = \qty{1.2e-4}{W/m^2}$,
$E_0 = \sqrt{2\Pi/\varepsilon_0c} = \qty{0.30}{V/m}$. (c) \qty{0.30}{V}. (d)
De retourstromen lopen door de grond; een weerstandbiedende bodem voegt
verliezen toe en bederft de spiegel; ingegraven radiale draden geven hun een
koperen weg.
\end{solution}

\begin{solution}{exo:b2:dipole-radiation:12}
(a) $p_0 = \alpha E_0$ geeft
$\sigma = \alpha^2\omega^4/6\pi\varepsilon_0^2c^4$; met
$\alpha = \varepsilon_0(n^2
- 1)/N$ en $\omega = 2\pi c/\lambda$:
$\sigma = 8\pi^3(n^2 - 1)^2/3N^2\lambda^4$. (b)
$n^2 - 1 =
5.8 \times 10^{-4}$: $\sigma = \qty{4.9e-31}{m^2}$, gemiddelde
vrije weglengte $1/N\sigma = \qty{80}{km}$. (c) $8/80 = 0.1$. (d) $\sigma$
wordt uit de \omterm{def:b2:dispersion-wave-packets:complex}{verzwakking} van de hemel gemeten en $n$ is bekend: daaruit
volgt $N$, en dus het getal van Avogadro.
\end{solution}

\begin{solution}{pb:b2:dipole-radiation:1}
\textbf{1.} $\lambda = \qty{300}{m}$, $h = \qty{75}{m}$;
$I_0 = \sqrt{2 \times 10^5/36} = \qty{75}{A}$; $V = RI_0 = \qty{2.7}{kV}$.

\textbf{2.} $\ell_{\text{eff}} = \qty{48}{m}$:
$p_0 = 75 \times 48/6.3 \times 10^6 = \qty{5.7e-4}{C.m}$.

\textbf{3.} $\langle\Pi\rangle = 1.5 \times 10^5/2\pi r^2$:
\qty{2.4e-4}{W/m^2} en \qty{2.4e-6}{W/m^2}; $E_0 = \qty{0.42}{V/m}$ en
\qty{0.042}{V/m}.

\textbf{4.} Spriet: \qty{42}{mV}. Ferrietstaaf:
$N\mu_{\text{eff}}A\omega B_0 = 10^4 \times 10^{-4} \times 6.3 \times
10^6 \times 1.4 \times 10^{-10} = \qty{0.9}{mV}$
--- minder, maar compact en selectief.

\textbf{5.} $A = 1.5\lambda^2/4\pi = \qty{1.1e4}{m^2}$: $P = \Pi A = \qty{26}{mW}$.

\textbf{6.} $36/38 = 95\%$; \qty{5}{kW} warmte.

\textbf{7.} $r/\lambda = 33$ op \qty{10}{km}: stralingszone; op \qty{100}{m}
is $r \approx \lambda/3$: het overgangsgebied, waar het quasistatische veld
nog meetelt.

\textbf{8.} Overdag absorbeert de D-laag de golf die anders de weerkaatsende
lagen zou bereiken; 's nachts is zij verdwenen en draagt de ruimtegolf de
zender honderden kilometers ver.

\textbf{9.} $\sigma(450) = \qty{1.0e-30}{m^2}$,
$\sigma(650) = \qty{2.4e-31}{m^2}$: verhouding $4.4$.

\textbf{10.} Gemiddelde vrije weglengten $39$, $87$ en \qty{170}{km};
$\tau = 0.21$, $0.09$, $0.05$: $19\%$, $9\%$ en $5\%$ verstrooid.

\textbf{11.}
$N\sigma I = 2.5 \times 10^{25} \times 4.6 \times 10^{-31} \times 300 = \qty{3.5}{mW/m^3}$,
in $4\pi$ met het diagram $(1 + \cos^2\chi)$ voor \omterm{def:b2:plane-waves-polarization:states}{natuurlijk licht}.

\textbf{12.} De kolom verstrooit
$3.5 \times 10^{-3} \times 8000 = \qty{28}{W}$ van haar \qty{300}{W}
($9\%$); uitgesmeerd over de halve bol levert de hemel zo'n
\qty{4}{W.m^{-2}.sr^{-1}}, tegen de
$300/6.8 \times 10^{-5} = \qty{4e6}{W.m^{-2}.sr^{-1}}$ van de zon: een
miljoenste per steradiaal --- de juiste orde.

\textbf{13.} Bij enkelvoudige verstrooiing volledig gepolariseerd op
\qty{90}{\degree}; meervoudige verstrooiing en de grond brengen dat tot zo'n
$75\%$ terug; bijen en mieren lezen het patroon om de zon achter de wolken
te vinden.

\textbf{14.} $\tau \times 38$: blauw $8$ ($T = 3 \times 10^{-4}$), rood
$1.8$ ($T = 0.17$): een diep rode zon; recht boven komt het licht na een
korte weg aan en is het blauw.

\textbf{15.} Geen dampkring, niets om te verstrooien: de hemel is zwart
naast de zon.

\textbf{16.} Stof van een micrometer verstrooit als miedeeltjes: vooruit en
rood, wat een botergele hemel geeft; vlak bij de ondergaande zon is het
vooruit verstrooide licht blauw.

\textbf{17.} $8\pi r_e^2/3 = \qty{6.65e-29}{m^2}$.

\textbf{18.} $a = v^2/2d = \qty{6.4e21}{m/s^2}$;
$\mathcal P = e^2a^2/6\pi\varepsilon_0c^3 =
\qty{2e-10}{W}$ gedurende
$2d/v = \qty{2.5e-14}{s}$: \qty{6e-24}{J}, $10^{-10}$ van \qty{100}{keV} ---
de werkelijke opbrengst ligt bij $1\%$, omdat het afremmen in heftige nabije
ontmoetingen met kernen gebeurt en niet in een zachte eenparige vertraging.

\textbf{19.} $hc/E = \qty{12.4}{pm}$.

\textbf{20.} \qty{2}{kW} in de bundel, ongeveer \qty{20}{W} röntgenstraling;
twee kilowatt warmte op één vlek --- de anode draait om die te spreiden.

\textbf{21.} $a = v^2/r = \qty{9e22}{m/s^2}$, $\mathcal P = \qty{5e-8}{W}$;
\qty{13.6}{eV} weg in \qty{5e-11}{s}.

\textbf{22.} Klassieke ladingen in gebonden beweging stralen onafgebroken;
de kwantummechanica heeft stationaire toestanden die dat niet doen, en een
laagste toestand waar niets onder ligt.

\textbf{23.} $a = c^2/R = \qty{9e14}{m/s^2}$: \qty{5e-24}{W} volgens Larmor,
in werkelijkheid $\gamma^4 \sim 10^{16}$ keer meer: tientallen kilowatt
synchrotronlicht uit een opgeslagen bundel.

\textbf{24.} $\sigma_TI = \qty{6.7e-26}{W}$; $1.5 \times 10^{25}$ elektronen
voor een watt.

\textbf{25.} Mast: \qty{5.7e-4}{C.m} bij \qty{1}{MHz}, \qty{100}{kW}, een
verticale torus. Molecuul: $\sim$\qty{1e-35}{C.m} bij \qty{500}{THz},
\qty{1e-28}{W}, een torus om het veld. Vrij elektron: Thomson,
\qty{1e-25}{W} in zonlicht, hetzelfde diagram.
\end{solution}
